\section{杨--米尔斯梯度流与小流时间展开}
\label{sec:3}
杨--米尔斯梯度流沿着一个虚构的时间~$t$，按照流方程
\begin{equation}
   \partial_tB_\mu(t,x)=D_\nu G_{\nu\mu}(t,x)
   +\alpha_0D_\mu\partial_\nu B_\nu(t,x),
\label{eq:(3.1)}
\end{equation}
定义一个 $D+1$~维的规范势~$B(t,x)$，其中 $D+1$~维场强与协变导数分别定义为
\begin{equation}
   G_{\mu\nu}(t,x)
   =\partial_\mu B_\nu(t,x)-\partial_\nu B_\mu(t,x)
   +[B_\mu(t,x),B_\nu(t,x)]
\label{eq:(3.2)}
\end{equation}
与
\begin{equation}
   D_\mu=\partial_\mu+[B_\mu,\cdot].
\label{eq:(3.3)}
\end{equation}
流的初始条件由前一节的 $D$~维规范势给出：
\begin{equation}
   B_\mu(t=0,x)=A_\mu(x).
\label{eq:(3.4)}
\end{equation}
在式~\eqref{eq:(3.1)} 中，引入最后一项是为了抑制场沿规范自由度方向的演
化。虽然这一项破坏规范对称性，但它不影响任何规范不变算符的演
化~\cite{Luscher:2010iy}。注意，流时间~$t$ 的质量量纲为~$-2$。

现在，从延拓到 $D+1$~维的场强~\eqref{eq:(3.2)} 出发，我们定义能量--动量张
量的一个 $D+1$~维类似物：
\begin{equation}
   U_{\mu\nu}(t,x)\equiv
   G_{\mu\rho}^a(t,x)G_{\nu\rho}^a(t,x)
   -\frac{1}{4}\delta_{\mu\nu}G_{\rho\sigma}^a(t,x)G_{\rho\sigma}^a(t,x).
\label{eq:(3.5)}
\end{equation}
虽然它在形式上类似于原来的能量--动量张量~\eqref{eq:(2.4)}，但这个
$D+1$~维对象与式~\eqref{eq:(2.4)} 之间如何关联（或者是否关联）并非先验显
然。找出二者之间的关系正是本文的主要任务。我们还使用文
献~\cite{Luscher:2010iy} 中研究的密度算符：
\begin{equation}
   E(t,x)\equiv\frac{1}{4}G_{\mu\nu}^a(t,x)G_{\mu\nu}^a(t,x).
\label{eq:(3.6)}
\end{equation}

如文献~\cite{Luscher:2011bx} 所示，对~$t>0$，$B_\mu(t,x)$ 的任何关联函数在
$D$~维杨--米尔斯理论的标准重正化之后都是 UV 有限的。这一性质甚至对
$B_\mu(t,x)$ 的任何局域乘积（例如式~\eqref{eq:(3.5)} 与~\eqref{eq:(3.6)}）
都成立。此外，在小流时间下，$B_\mu(t,x)$ 的局域乘积可以在 $D$~维的意义上
视为局域场，因为流方程~\eqref{eq:(3.1)} 基本上是沿时间~$t$ 的扩散方程，而
$x$ 方向上的扩散长度为~$\sqrt{8t}$。这些性质使我们能够像文献%
~\cite{Luscher:2011bx} 第 8 节解释的那样，把 $U_{\mu\nu}(t,x)$
与~$E(t,x)$ 表示成带有限系数的 $D$~维重正化局域算符的渐近级数。考虑到规范
不变性与指标结构，对 $D=4$ 可以写成
\begin{equation}
   U_{\mu\nu}(t,x)
   =c_T(t)\left\{T_{\mu\nu}\right\}_R(x)
   +c_S(t)\delta_{\mu\nu}
   \left\{\frac{1}{4}F_{\rho\sigma}^aF_{\rho\sigma}^a\right\}_R(x)
   +O(t),
\label{eq:(3.7)}
\end{equation}
其中被省略的项是质量量纲大于等于~$6$ 的算符的贡献。对于式~\eqref{eq:(3.6)}，
我们类似地有
\begin{equation}
   E(t,x)
   =\left\langle E(t,x)\right\rangle
   +c_E(t)\left\{\frac{1}{4}F_{\rho\sigma}^aF_{\rho\sigma}^a\right\}_R(x)
   +O(t).
\label{eq:(3.8)}
\end{equation}
我们注意到，固定重正化规范耦合时，$U_{\mu\nu}(t,x)$~\eqref{eq:(3.5)} 在
$D=4$ 时是无迹的：
\begin{equation}
   \delta_{\mu\nu}U_{\mu\nu}(t,x)
   =2\epsilon E(t,x)\xrightarrow{\epsilon\to0}0,
\label{eq:(3.9)}
\end{equation}
因为 $E(t,x)$~\eqref{eq:(3.6)} 是有限的~\cite{Luscher:2010iy} 且不产生
$1/\epsilon$ 奇异性（这解释了为什么式~\eqref{eq:(3.7)} 中没有 $c$ 数期望值
项）。于是，考虑式~\eqref{eq:(3.7)} 的迹部分可以看到：由于迹反常%
~\eqref{eq:(2.15)}，系数 $c_T(t)$ 与~$c_S(t)$ 并不独立，在 $D=4$
 时它们由
\begin{equation}
   c_S(t)=\frac{\beta}{2g^3}c_T(t)
\label{eq:(3.10)}
\end{equation}
相联系。

从式~\eqref{eq:(3.7)} 与~\eqref{eq:(3.8)} 中消去重正化的作用量密度，得到
\begin{equation}
   \left\{T_{\mu\nu}\right\}_R(x)
   =\frac{1}{c_T(t)}U_{\mu\nu}(t,x)
   -\frac{c_S(t)}{c_T(t)c_E(t)}\delta_{\mu\nu}
   \left[E(t,x)
   -\left\langle E(t,x)\right\rangle
   \right]
   +O(t).
\label{eq:(3.11)}
\end{equation}
这个表达式把能量--动量张量~\eqref{eq:(2.6)} 与梯度流定义的规范不变局域乘积
的短流时间行为联系起来。因此，一旦知道了这些系数，就可以从右边组合的
$t\to0$ 行为中提取能量--动量张量。
